% Chapter 1
\chapter{Definición del problema u oportunidad} % Main chapter title
\label{sec:problema} % For referencing the chapter elsewhere, use

\section{Problema}
El problema central de esta investigación radica en la severa degradación de la adherencia a la secuencia de clientes planificada por la heurística de SimpliRoute. Así, en el contexto de la última milla, el software modela la operación como un Problema de Ruteo de Vehículos (VRP)[comentario:no se si el software lo modela directamente como una vrp, no se nada del software, solo se que es una heuristica], donde el objetivo principal es determinar una permutación ordenada de paradas para cada vehículo, minimizando una función de costo basada en distancias y tiempos teóricos de viaje. Sin embargo, los datos empíricos de la plataforma revelan una desconexión crítica con la realidad. Ya que, la tasa de adherencia de los conductores a la secuencia exacta de clientes propuesta por Simpliroute alcanza apenas un 14.19\% \cite{rubiocarrasco2025mejora}. Asimismo, esto significa que en el 85.81\% de los casos, los operadores alteran deliberadamente el orden de visita de los nodos o clientes a visitar. De manera que, este fenómeno representa un desajuste estructural entre el espacio de soluciones del algoritmo y las decisiones secuenciales óptimas del operador humano.

La persistencia de esta baja adherencia impacta negativamente en la eficiencia proyectada por la plataforma, alterando las estimaciones de tiempos de llegada y de los costos logísticos, arriesgando la confiabilidad del servicio de SimpliRoute.

\section{Análisis de las Causas Basado en la Literatura}

La drástica desviación conductual de los operadores logísticos responde a variables críticas que la heurística tradicional omite. Este desajuste puede descomponerse en tres lógicas principales para entender el problema.

\subsection{Insuficiencia del criterio único de optimización}

Primero, Los motores de ruteo tradicionales resuelven el Problema de Ruteo de Vehículos ($VRP$) modelando una función de costo lineal que minimiza variables geométricas o temporales explícitas, tales como la distancia total o el tiempo teórico de viaje. Sin embargo, la evidencia empírica señala que los conductores expertos operan bajo un esquema multiobjetivo subjetivo y rara vez eligen el camino más corto o rápido de forma aislada 

La literatura identifica que los operadores humanos evalúan de forma latente la calidad de la infraestructura y el estrés de la conducción. Por ejemplo, desvían sus rutas deliberadamente para evitar segmentos viales con pavimentos en mal estado, zonas con exceso de semáforos, giros complejos a la izquierda en avenidas saturadas, entre otros \cite{zattoni2025inverse} \cite{cook2024constrained}\cite{yang2015toward} \cite{pan2025preference}\cite{ceikute2013routing} \cite{chen2001using}. Al omitir estas variables en la función de costo del planificador, el espacio de soluciones del algoritmo se vuelve impracticable en la realidad.

\subsection{Conocimiento tácito y preferencias contextuales}
Segundo, la causa crítica es la exclusión del conocimiento asimétrico que posee el conductor frente a la plataforma logística. A través de la repetición sistemática de sus rutas, los repartidores desarrollan una profunda familiaridad con sus zonas asignadas, acumulando un "conocimiento tácito" (tacit knowledge) que no está digitalizado en las bases de datos del sistema \cite{merchan20242021} \cite{canoy2023learn} \cite{ceikute2013routing} \cite{ozarik2024machine}.

Este aprendizaje basado en la experiencia les permite anticipar dinámicas locales que el algoritmo ignora de forma nativa. Entre ellas destacan la disponibilidad real de estacionamiento en zonas de alta densidad, las ventanas efectivas de recepción de los clientes (que suelen diferir de las ventanas teóricas ingresadas en la plataforma) y factores del entorno como los horarios de congestión inducida por escuelas o comercios locales (Rastogi, 2023; Ulmer y cols., 2020). Al verse enfrentado a una secuencia algorítmica que contradice esta experiencia, el conductor prioriza su propio criterio heurístico para garantizar la efectividad de la entrega (Li y Phillips, 2018).

\subsection{Rigidez de los modelos estáticos ante contextos dinámicos}

Por último, la literatura atribuye las desviaciones a la naturaleza determinista o estática con la que se suelen calcular las matrices de tiempo de viaje en los planificadores comerciales. Mientras que la plataforma calcula un plan rígido al inicio de la jornada utilizando velocidades promedio, el entorno de la última milla es altamente estocástico y sensible al contexto (Guo y cols., 2020).

Los cambios en las condiciones climáticas, la congestión vehicular imprevista o eventos disruptivos en la red vial obligan al conductor a tomar decisiones adaptativas en tiempo real basándose en información local (Tang y Levinson, 2018). Al no existir un mecanismo dinámico que permita al motor de optimización aprender cómo varían las preferencias de ruta según el contexto horario o espacial de cada conductor individual, el plan original se vuelve obsoleto a las pocas horas de iniciar la operación.


\section{Pregunta de Investigación}
En base a los antecedentes expuestos, la presente investigación busca responder a la siguiente interrogante ¿Cómo se pueden modelar e integrar las preferencias operacionales latentes de los conductores, inferidas a partir de las discrepancias entre rutas planificadas y trazas GPS históricas, para lograr una mayor adherencia del motor de optimización vehicular?